\documentclass[11pt]{article}
\usepackage[utf-8]{inputenc}
\usepackage[spanish]{babel}
\usepackage{hyperref}
\usepackage{xcolor}
\usepackage{array}
\usepackage{booktabs}

\title{Acceso a Plataforma Staging --- MultiCatálogo}
\author{Andy Alejandro}
\date{\today}

\begin{document}

\maketitle

\section{URL Pública}

\begin{center}
\Large\textbf{\textcolor{blue}{\url{https://staging-andy.tailc8cd10.ts.net/}}}
\end{center}

\textbf{Requisito:} Tailscale instalado y conectado al tailnet de la UPSE.

\section{Credenciales de Usuarios}

\subsection{Administrador}

\begin{table}[h]
\centering
\begin{tabular}{|l|l|}
\hline
\textbf{Campo} & \textbf{Valor} \\
\hline
Email & \texttt{admin@upse.edu.ec} \\
\hline
Contraseña & \texttt{12345678} \\
\hline
Rol & Admin \\
\hline
Nombre & Administrador \\
\hline
\end{tabular}
\end{table}

\textbf{Permisos:}
\begin{itemize}
\item Panel de administración
\item Gestión de productos
\item Reportes y auditoría
\item Auditoría de comisiones
\end{itemize}

\subsection{Cliente Demo}

\begin{table}[h]
\centering
\begin{tabular}{|l|l|}
\hline
\textbf{Campo} & \textbf{Valor} \\
\hline
Email & \texttt{cliente@upse.edu.ec} \\
\hline
Contraseña & \texttt{12345678} \\
\hline
Rol & Cliente \\
\hline
Nombre & Cliente Demo \\
\hline
\end{tabular}
\end{table}

\textbf{Permisos:}
\begin{itemize}
\item Visualizar catálogo de productos
\item Realizar compras
\item Gestionar red de referidos
\item Consultar comisiones
\end{itemize}

\section{Instrucciones de Acceso}

\begin{enumerate}
\item Instala Tailscale en tu dispositivo (disponible para Mac, Windows, Linux, iOS, Android)
\item Conecta al tailnet de la UPSE
\item Accede a: \url{https://staging-andy.tailc8cd10.ts.net/}
\item Inicia sesión con cualquiera de las credenciales anteriores
\end{enumerate}

\section{Arquitectura del Sistema}

\begin{table}[h]
\centering
\begin{tabular}{|l|l|l|}
\hline
\textbf{Componente} & \textbf{URL} & \textbf{Puerto} \\
\hline
Frontend & \url{https://staging-andy.tailc8cd10.ts.net/} & 5173 \\
\hline
Backend API & \url{https://staging-andy.tailc8cd10.ts.net/:3000} & 3000 \\
\hline
Health Check & \url{https://staging-andy.tailc8cd10.ts.net/:3000/api/health} & --- \\
\hline
\end{tabular}
\end{table}

\section{Notas de Desarrollo}

\begin{itemize}
\item \textbf{Base de datos:} PostgreSQL 16 (puerto 5433 local)
\item \textbf{Caché:} Redis (puerto 6379 local)
\item \textbf{Contraseña seed:} Válida solo para desarrollo
\item \textbf{Archivo seed:} \texttt{Talleres/Practica07/database/migrations/003\_seed.sql}
\item \textbf{Despliegue:} Automático en push a rama \texttt{develop}
\end{itemize}

\section{Servidor de Despliegue}

\begin{table}[h]
\centering
\begin{tabular}{|l|l|}
\hline
\textbf{Característica} & \textbf{Valor} \\
\hline
Hostname & staging-andy \\
\hline
IP Interna & 10.10.10.26 \\
\hline
Ubicación & Proxmox VM 109 \\
\hline
Sistema Operativo & Debian \\
\hline
Motor de Contenedores & Docker Compose \\
\hline
\end{tabular}
\end{table}

\section{Troubleshooting}

\subsection{No puedo acceder desde mi dispositivo}

\begin{enumerate}
\item Verifica que Tailscale esté \textbf{activo} en tu dispositivo
\item Verifica que estés en el \textbf{mismo tailnet} que la VM
\item Intenta: \texttt{ping staging-andy.tailc8cd10.ts.net} desde terminal
\end{enumerate}

\subsection{El backend responde 404}

\begin{enumerate}
\item Verifica contenedores: \texttt{docker compose ps}
\item Revisa logs: \texttt{docker compose logs backend | tail -20}
\end{enumerate}

\subsection{El frontend no carga las rutas (SPA)}

\begin{enumerate}
\item Verifica serve: \texttt{docker compose logs frontend}
\item Reconstruye: \texttt{docker compose down -v \&\& docker compose --profile staging up -d --build}
\end{enumerate}

\vspace{2cm}

\begin{center}
\textbf{Última actualización:} 23 de octubre de 2026 \\
\textbf{Ambiente:} Staging (Tailscale Funnel) \\
\textbf{Responsable:} Andy Alejandro
\end{center}

\end{document}
